\chapter{Chapter09 — Hypercontractive Applications}
%%% generated by the blueprint pipeline from FABL.Chapter09.HypercontractiveApplications
\section{Overview}

Declarations of the Lean module \texttt{FABL.Chapter09.HypercontractiveApplications}.

\section{Declarations}

\begin{theorem}[Power of the uniform norm]
\lean{FABL.uniformLpNorm_rpow}
\label{FABL.uniformLpNorm_rpow}
\leanok
\uses{FABL.uniformLpNorm}
\difficulty{3}
Let $\Omega$ be a finite set, let $q>0$, and let $f\colon\Omega\to\mathbb R$. The $q$th
power of the uniform $L^q$ norm of $f$ equals the uniform average of $|f|^q$.
\end{theorem}

\begin{theorem}[Power scaling of the uniform norm]
\lean{FABL.uniformLpNorm_abs_rpow}
\label{FABL.uniformLpNorm_abs_rpow}
\leanok
\uses{FABL.uniformLpNorm}
\difficulty{4}
For a finite set $\Omega$, real numbers $p,c>0$, and a function $f:\Omega\to\mathbb{R}$,
the uniform $L^{p/c}$ quantity of $x\mapsto |f(x)|^c$ equals the uniform $L^p$ quantity
of $f$, raised to the power $c$.
\end{theorem}

\begin{theorem}[Interpolation inequality for the uniform norm]
\lean{FABL.uniformLpNorm_two_interpolation_sq}
\label{FABL.uniformLpNorm_two_interpolation_sq}
\leanok
\uses{FABL.uniformInner, FABL.uniformInner_le_uniformLpNorm_mul, FABL.uniformLpNorm, FABL.uniformLpNorm_abs_rpow, FABL.uniformLpNorm_two_sq_eq_expect_sq}
\difficulty{5}
Let $\Omega$ be a finite nonempty set, let $p,q,\theta\in\mathbb R$ satisfy $p>0$,
$q>0$, $0<\theta<1$, and $\frac12=\frac{\theta}{p}+\frac{1-\theta}{q}$. For every
$f:\Omega\to\mathbb R$, the uniform norms satisfy $\|f\|_2^2\le
\|f\|_p^{2\theta}\|f\|_q^{2(1-\theta)}$.
\end{theorem}

\begin{theorem}[Low-degree norm interpolation bound]
\lean{FABL.lowDegree_two_le_interpolation_constant}
\proofsource{Analysis-of-Boolean-Functions-by-Ryan-ODonnell}{
  pdf-ffbe9c7b3306-p262-b006,
  pdf-ffbe9c7b3306-p273-b002
}
\label{FABL.lowDegree_two_le_interpolation_constant}
\leanok
\uses{FABL.fourierDegree, FABL.lowDegree_uniformLpNorm_le, FABL.uniformLpNorm, FABL.uniformLpNorm_nonneg, FABL.uniformLpNorm_two_interpolation_sq}
\difficulty{6}
Let $f$ be a real-valued function on the $n$-dimensional sign cube whose Fourier degree
is at most $k$. For $1\le p<2$ and $\varepsilon>0$, its uniform $L^2$ norm satisfies
\[\|f\|_2\le (1+\varepsilon)^{(2/p-1)k(1/\varepsilon+1/2)}\|f\|_p.\]
\end{theorem}

\begin{theorem}[Interpolation constant exponential bound]
\lean{FABL.interpolation_constant_le_exp}
\label{FABL.interpolation_constant_le_exp}
\leanok
\difficulty{3}
Let $c,psilon\in\mathbb{R}$ satisfy $c\ge 0$ and $\varepsilon>0$. Then
\[
(1+\varepsilon)^{c(1/\varepsilon+1/2)}\le \exp\bigl(c(1+\varepsilon/2)\bigr).
\]
\end{theorem}

\begin{theorem}[Comparison of low degree norms]
\lean{FABL.lowDegree_two_le_expFactor_mul}
\proofsource{Analysis-of-Boolean-Functions-by-Ryan-ODonnell}{pdf-ffbe9c7b3306-p262-b006}
\label{FABL.lowDegree_two_le_expFactor_mul}
\leanok
\uses{FABL.fourierDegree, FABL.interpolation_constant_le_exp, FABL.lowDegree_two_le_interpolation_constant, FABL.uniformLpNorm, FABL.uniformLpNorm_nonneg}
\difficulty{6}
Let $f$ be a real-valued function on the $n$-dimensional Boolean cube with Fourier
degree at most $k$, where $k\in\mathbb N$. For $1\le p\le 2$, its uniform norms satisfy
$\|f\|_2\le e^{k(2/p-1)}\|f\|_p$.
\end{theorem}

\begin{theorem}[Uniform norm bound for low degree functions]
\lean{FABL.lowDegree_two_le_exp_pow_mul_one}
\proofsource{Analysis-of-Boolean-Functions-by-Ryan-ODonnell}{pdf-ffbe9c7b3306-p262-b006}
\label{FABL.lowDegree_two_le_exp_pow_mul_one}
\leanok
\uses{FABL.fourierDegree, FABL.lowDegree_two_le_expFactor_mul, FABL.uniformLpNorm}
\difficulty{3}
Let $f$ be a real-valued function on the $n$-dimensional Boolean cube with Fourier
degree at most $k$, where $k\in\mathbb N$. Its uniform norms satisfy $\|f\|_2\le
e^k\|f\|_1$.
\end{theorem}

\begin{theorem}[Monotonicity of the uniform norm]
\lean{FABL.uniformLpNorm_mono}
\proofsource{Analysis-of-Boolean-Functions-by-Ryan-ODonnell}{pdf-ffbe9c7b3306-p037-b002}
\label{FABL.uniformLpNorm_mono}
\leanok
\uses{FABL.abs_expect_le_uniformLpNorm, FABL.uniformLpNorm, FABL.uniformLpNorm_abs_rpow, FABL.uniformLpNorm_nonneg, FABL.uniformLpNorm_rpow}
\difficulty{5}
Let $\Omega$ be a finite nonempty set, let $0<p\le q$, and let $f\colon\Omega\to\mathbb
R$. Then the uniform $L^p$ norm of $f$ is at most its uniform $L^q$ norm.
\end{theorem}

\begin{theorem}[Degree of a linear form]
\lean{FABL.fourierDegree_linearForm_le_one}
\label{FABL.fourierDegree_linearForm_le_one}
\leanok
\uses{FABL.expect_monomial_mul, FABL.fourierCoeff, FABL.fourierDegree, FABL.fourierDegree_le_iff, FABL.linearForm, FABL.monomial, FABL.signValue}
\difficulty{5}
For every real coefficient vector $a=(a_i)_{i=1}^n$, the Fourier degree of the linear
form $x\\mapsto\\sum_{i=1}^n a_i x_i$ on the Boolean cube is at most $1$.
\end{theorem}

\begin{theorem}[Khintchine of le two]
\lean{FABL.khintchine_of_le_two}
\statementsource{Analysis-of-Boolean-Functions-by-Ryan-ODonnell}{
  pdf-ffbe9c7b3306-p115-b002,
  pdf-ffbe9c7b3306-p262-b005
}
\label{FABL.khintchine_of_le_two}
\leanok
\uses{FABL.fourierDegree_linearForm_le_one, FABL.linearForm, FABL.lowDegree_two_le_expFactor_mul, FABL.uniformLpNorm, FABL.uniformLpNorm_mono}
\difficulty{4}
% dataset-pipeline: stub -- prose pending build_informal_statements.py
Statement pending regeneration.
\end{theorem}

\begin{theorem}[Khintchine of two le]
\lean{FABL.khintchine_of_two_le}
\statementsource{Analysis-of-Boolean-Functions-by-Ryan-ODonnell}{pdf-ffbe9c7b3306-p262-b005}
\label{FABL.khintchine_of_two_le}
\leanok
\uses{FABL.fourierDegree_linearForm_le_one, FABL.linearForm, FABL.lowDegree_uniformLpNorm_le, FABL.uniformLpNorm, FABL.uniformLpNorm_mono}
\difficulty{3}
% dataset-pipeline: stub -- prose pending build_informal_statements.py
Statement pending regeneration.
\end{theorem}

\begin{theorem}[Khintchine inequality]
\lean{FABL.khintchineInequality}
\statementsource{Analysis-of-Boolean-Functions-by-Ryan-ODonnell}{pdf-ffbe9c7b3306-p262-b005}
\label{FABL.khintchineInequality}
\leanok
\uses{FABL.khintchine_of_le_two, FABL.khintchine_of_two_le, FABL.linearForm, FABL.uniformLpNorm, FABL.uniformLpNorm_nonneg}
\difficulty{5}
% dataset-pipeline: stub -- prose pending build_informal_statements.py
Statement pending regeneration.
\end{theorem}

\begin{theorem}[Fourier degree after mean subtraction]
\lean{FABL.fourierDegree_sub_mean_le}
\label{FABL.fourierDegree_sub_mean_le}
\leanok
\uses{FABL.expect_monomial, FABL.fourierCoeff, FABL.fourierDegree, FABL.fourierDegree_le_iff, FABL.mean, FABL.monomial}
\difficulty{4}
Let $f:\{-1,1\}^n\to\mathbb{R}$, and let $k$ be a natural number. If the Fourier degree
of $f$ is at most $k$, then the Fourier degree of $x\mapsto f(x)-\mathbb{E}[f]$, where
the expectation is uniform on the sign cube, is at most $k$.
\end{theorem}

\begin{theorem}[Probability of exceeding the mean]
\lean{FABL.lowDegree_aboveMean_probability}
\proofsource{Analysis-of-Boolean-Functions-by-Ryan-ODonnell}{
  pdf-ffbe9c7b3306-p263-b005,
  pdf-ffbe9c7b3306-p263-b006
}
\label{FABL.lowDegree_aboveMean_probability}
\leanok
\uses{FABL.fourierDegree, FABL.fourierDegree_sub_mean_le, FABL.lowDegree_two_le_exp_pow_mul_one, FABL.mean, FABL.pmfExpectation_mul_sq_le_sq_mul_sq, FABL.pmfExpectation_uniformPMF_eq_expect, FABL.uniformLpNorm, FABL.uniformLpNorm_nonneg, FABL.uniformLpNorm_two_sq_eq_expect_sq, FABL.uniformPMF, FABL.uniformProbability}
\difficulty{6}
Let $f$ be a real-valued function on the $n$-dimensional Boolean cube with Fourier
degree at most $k$, where $k\in\mathbb N$, and suppose $f$ is nonconstant. For uniformly
random $x$ in the cube, $\Pr[f(x)>\mathbb E f]\ge \frac14 e^{-2k}$.
\end{theorem}

\begin{theorem}[Norm of a zero-one-valued function]
\lean{FABL.uniformLpNorm_zeroOne}
\proofsource{Analysis-of-Boolean-Functions-by-Ryan-ODonnell}{pdf-ffbe9c7b3306-p024-b001}
\label{FABL.uniformLpNorm_zeroOne}
\leanok
\uses{FABL.mean, FABL.uniformLpNorm}
\difficulty{3}
Let $f$ be a real-valued function on the $n$-dimensional Boolean cube, and suppose every
value of $f$ is either $0$ or $1$. If its uniform mean is $\alpha$ and $p>0$, then its
uniform $L^p$ norm is $\alpha^{1/p}$.
\end{theorem}

\begin{theorem}[Noise stability le uniform lp norm one add sq]
\lean{FABL.noiseStability_le_uniformLpNorm_oneAdd_sq}
\statementsource{Analysis-of-Boolean-Functions-by-Ryan-ODonnell}{
  pdf-ffbe9c7b3306-p254-b002,
  pdf-ffbe9c7b3306-p254-b003,
  pdf-ffbe9c7b3306-p260-b001
}
\label{FABL.noiseStability_le_uniformLpNorm_oneAdd_sq}
\leanok
\uses{FABL.noiseOperator, FABL.noiseOperator_comp, FABL.noiseStability, FABL.noiseStability_eq_uniformInner_noiseOperator, FABL.pTwoHypercontractivity, FABL.uniformInner, FABL.uniformInner_noiseOperator_left, FABL.uniformLpNorm, FABL.uniformLpNorm_nonneg, FABL.uniformLpNorm_two_sq_eq_uniformInner}
\difficulty{4}
% dataset-pipeline: stub -- prose pending build_informal_statements.py
Statement pending regeneration.
\end{theorem}

\begin{theorem}[Small-set noise stability bound]
\lean{FABL.smallSetExpansion}
\proofsource{Analysis-of-Boolean-Functions-by-Ryan-ODonnell}{pdf-ffbe9c7b3306-p264-b001}
\label{FABL.smallSetExpansion}
\leanok
\uses{FABL.mean, FABL.noiseStability, FABL.noiseStability_le_uniformLpNorm_oneAdd_sq, FABL.uniformLpNorm, FABL.uniformLpNorm_zeroOne}
\difficulty{4}
Let $f:\{-1,1\}^n\to\mathbb{R}$ take only the values $0$ and $1$, and suppose its
uniform mean is $\alpha$. For $\rho\in[0,1]$, the noise stability of $f$ at correlation
$\rho$ is at most $\alpha^{2/(1+\rho)}$.
\end{theorem}

\begin{theorem}[Small-set expansion event bound]
\lean{FABL.smallSetExpansion_event}
\proofsource{Analysis-of-Boolean-Functions-by-Ryan-ODonnell}{
  pdf-ffbe9c7b3306-p264-b001,
  pdf-ffbe9c7b3306-p254-b002,
  pdf-ffbe9c7b3306-p254-b003,
  pdf-ffbe9c7b3306-p254-b005,
  pdf-ffbe9c7b3306-p254-b006
}
\label{FABL.smallSetExpansion_event}
\leanok
\uses{FABL.correlatedPairPMF, FABL.mean, FABL.noiseStability, FABL.pmfEventProbability, FABL.pmfExpectation, FABL.smallSetExpansion}
\difficulty{4}
Let $f:\{-1,1\}^n\to\mathbb{R}$ take only the values $0$ and $1$, and suppose its
uniform mean is $\alpha$. For $\rho\in[0,1]$, under the correlated-pair distribution at
correlation $\rho$, the probability that both coordinates of the pair lie in the set
where $f=1$ is at most $\alpha^{2/(1+\rho)}$.
\end{theorem}

\begin{theorem}[Small set expansion conditional]
\lean{FABL.smallSetExpansion_conditional}
\statementsource{Analysis-of-Boolean-Functions-by-Ryan-ODonnell}{
  pdf-ffbe9c7b3306-p264-b001,
  pdf-ffbe9c7b3306-p254-b002,
  pdf-ffbe9c7b3306-p254-b003,
  pdf-ffbe9c7b3306-p254-b005,
  pdf-ffbe9c7b3306-p254-b006
}
\label{FABL.smallSetExpansion_conditional}
\leanok
\uses{FABL.correlatedPairPMF, FABL.mean, FABL.pmfEventProbability, FABL.smallSetExpansion_event}
\difficulty{3}
% dataset-pipeline: stub -- prose pending build_informal_statements.py
Statement pending regeneration.
\end{theorem}

\begin{theorem}[Collision bound for unbiased labels]
\lean{FABL.unbiasedLabel_collisionProbability}
\label{FABL.unbiasedLabel_collisionProbability}
\leanok
\uses{FABL.correlatedPairPMF, FABL.mean, FABL.pmfEventProbability, FABL.pmfEventProbability_congr, FABL.pmfExpectation, FABL.pmfExpectation_sum, FABL.smallSetExpansion_event, FABL.uniformProbability}
\difficulty{6}
Let $n$ be a natural number, let $m$ be a positive integer, and let
$f:\{-1,1\}^n\to\{0,\ldots,m-1\}$ assign labels to cube points. Suppose each label is
attained by exactly a fraction $1/m$ of the cube. Choose $X$ uniformly from the cube
and, independently in each coordinate, form $Y$ by keeping $X$'s coordinate with
probability $(1+\rho)/2$ and flipping it otherwise, where $0\le\rho\le1$. Then
$\Pr[f(X)=f(Y)]\le (1/m)^{(1-\rho)/(1+\rho)}$.
\end{theorem}

\begin{theorem}[Derivative and pivotal indicator]
\lean{FABL.abs_discreteDerivative_toReal_eq_pivotalIndicator}
\proofsource{Analysis-of-Boolean-Functions-by-Ryan-ODonnell}{pdf-ffbe9c7b3306-p047-b003}
\label{FABL.abs_discreteDerivative_toReal_eq_pivotalIndicator}
\leanok
\uses{FABL.BooleanFunction, FABL.BooleanFunction.toReal, FABL.IsPivotal, FABL.discreteDerivative, FABL.pivotalIndicator, FABL.sq_discreteDerivative_toReal_eq_pivotalIndicator}
\difficulty{4}
Let $f$ be a Boolean function on the sign cube $\\{−1,1\\}^n$, let $i$ be a coordinate,
and let $x$ be a point of the cube. The absolute value of the discrete derivative of the
real-valued version of $f$ at $x$ in coordinate $i$ equals the indicator that flipping
coordinate $i$ changes $f(x)$.
\end{theorem}

\begin{theorem}[Stable influence bound by Boolean influence]
\lean{FABL.stableInfluence_le_booleanInfluence_rpow}
\proofsource{Analysis-of-Boolean-Functions-by-Ryan-ODonnell}{pdf-ffbe9c7b3306-p264-b003}
\label{FABL.stableInfluence_le_booleanInfluence_rpow}
\leanok
\uses{FABL.BooleanFunction, FABL.BooleanFunction.toReal, FABL.abs_discreteDerivative_toReal_eq_pivotalIndicator, FABL.booleanInfluence, FABL.booleanInfluence_eq_influence_toReal, FABL.discreteDerivative, FABL.influence, FABL.mean, FABL.noiseStability, FABL.pivotalIndicator, FABL.pmfExpectation_mono, FABL.smallSetExpansion, FABL.sq_discreteDerivative_toReal_eq_pivotalIndicator, FABL.stableInfluence, FABL.stableInfluence_eq_noiseStability_discreteDerivative}
\difficulty{4}
Let $n$ be a natural number, let $f:\{−1,1\}^n\to\{−1,1\}$, and let $i$ be a coordinate.
For $0\leq\rho\leq1$, the stable influence of the real-valued version of $f$ at $i$ is
at most the Boolean influence of $f$ at $i$ raised to the power $2/(1+\rho)$.
\end{theorem}

\begin{theorem}[Level-One Fourier Weight and Variance]
\lean{FABL.fourierWeightAtLevel_one_le_variance}
\proofsource{Analysis-of-Boolean-Functions-by-Ryan-ODonnell}{
  pdf-ffbe9c7b3306-p027-b001,
  pdf-ffbe9c7b3306-p028-b002,
  pdf-ffbe9c7b3306-p126-b001
}
\label{FABL.fourierWeightAtLevel_one_le_variance}
\leanok
\uses{FABL.fourierCoeff, FABL.fourierWeightAtLevel, FABL.variance, FABL.variance_eq_sum_sq_fourierCoeff}
\difficulty{3}
For any function $f$ from the $n$-dimensional sign cube to $\bbr$, the sum of the squared
Fourier coefficients indexed by singleton subsets is at most the variance of $f$ under
the uniform distribution on the cube.
\end{theorem}

\begin{theorem}[Level one inequality sharp]
\lean{FABL.levelOneInequality_sharp}
\statementsource{Analysis-of-Boolean-Functions-by-Ryan-ODonnell}{
  pdf-ffbe9c7b3306-p126-b002,
  pdf-ffbe9c7b3306-p126-b003,
  pdf-ffbe9c7b3306-p273-b005
}
\label{FABL.levelOneInequality_sharp}
\leanok
\uses{FABL.fourierWeightAtLevel, FABL.fourierWeightAtLevel_one_le_variance, FABL.mean, FABL.sharpLevelOneInequality, FABL.sharpLevelOneInequality_eq_zero, FABL.variance, FABL.variance_eq_sum_sq_fourierCoeff}
\difficulty{5}
% dataset-pipeline: stub -- prose pending build_informal_statements.py
Statement pending regeneration.
\end{theorem}

\begin{theorem}[Low-degree weight and noise stability]
\lean{FABL.rho_pow_mul_fourierWeightAtMost_le_noiseStability}
\proofsource{Analysis-of-Boolean-Functions-by-Ryan-ODonnell}{
  pdf-ffbe9c7b3306-p056-b002,
  pdf-ffbe9c7b3306-p056-b003,
  pdf-ffbe9c7b3306-p264-b006
}
\label{FABL.rho_pow_mul_fourierWeightAtMost_le_noiseStability}
\leanok
\uses{FABL.fourierCoeff, FABL.fourierWeight, FABL.fourierWeightAtMost, FABL.noiseStability, FABL.noiseStability_eq_sum_rho_pow_mul_sq_fourierCoeff}
\difficulty{3}
Let $n$ and $k$ be natural numbers, let $f:\{−1,1\}^n\to\mathbb R$, and let
$0<\rho\le1$. Write $W_{\le k}(f)=\sum_{|S|\le k}\widehat f(S)^2$, where $\widehat f(S)$
is the uniform Fourier coefficient, and let $\operatorname{Stab}_\rho(f)=\mathbb
E[f(X)f(Y)]$ for a uniform $X$ and a $Y$ obtained by independently retaining each
coordinate of $X$ with probability $(1+\rho)/2$ and flipping it otherwise. Then $\rho^k
W_{\le k}(f)\le\operatorname{Stab}_\rho(f)$.
\end{theorem}

\begin{theorem}[Low-Degree Fourier Weight Bound]
\lean{FABL.fourierWeightAtMost_le_invPow_mul_rpow}
\proofsource{Analysis-of-Boolean-Functions-by-Ryan-ODonnell}{pdf-ffbe9c7b3306-p264-b006}
\label{FABL.fourierWeightAtMost_le_invPow_mul_rpow}
\leanok
\uses{FABL.fourierWeightAtMost, FABL.mean, FABL.rho_pow_mul_fourierWeightAtMost_le_noiseStability, FABL.smallSetExpansion}
\difficulty{4}
Let $k$ be a natural number, let $f$ be a real-valued function on the $n$-dimensional
sign cube, and let $\alpha,\rho\in\bbr$. Suppose $f$ takes only the values $0$ and $1$,
its mean is $\alpha$, and $0<\rho\le 1$. Then the sum of its squared Fourier
coefficients over subsets of size at most $k$ is at most
$(\rho^{-1})^k\alpha^{2/(1+\rho)}$.
\end{theorem}

\begin{theorem}[Fourier weight up to degree k]
\lean{FABL.levelKInequality}
\proofsource{Analysis-of-Boolean-Functions-by-Ryan-ODonnell}{
  pdf-ffbe9c7b3306-p264-b005,
  pdf-ffbe9c7b3306-p264-b006
}
\label{FABL.levelKInequality}
\leanok
\uses{FABL.fourierWeightAtMost, FABL.fourierWeightAtMost_le_invPow_mul_rpow, FABL.mean}
\difficulty{6}
Let $f:\{-1,1\}^n\to\mathbb R$ take only the values $0$ and $1$, and let its mean be
$\alpha$. If $k$ is a positive integer satisfying $k\le 2\log(1/\alpha)$, then the sum
of the squared Fourier coefficients of $f$ over all subsets of coordinates of size at
most $k$ is at most $\left(\frac{2e}{k}\log\frac{1}{\alpha}\right)^k\alpha^2$.
\end{theorem}

\begin{theorem}[Level one weight and noise stability]
\lean{FABL.rho_mul_fourierWeightAtLevel_one_le_stability_sub_mean_sq}
\proofsource{Analysis-of-Boolean-Functions-by-Ryan-ODonnell}{
  pdf-ffbe9c7b3306-p056-b002,
  pdf-ffbe9c7b3306-p056-b003
}
\label{FABL.rho_mul_fourierWeightAtLevel_one_le_stability_sub_mean_sq}
\leanok
\uses{FABL.fourierCoeff, FABL.fourierWeight, FABL.fourierWeightAtLevel, FABL.mean, FABL.mean_eq_fourierCoeff_empty, FABL.noiseStability, FABL.noiseStability_eq_sum_rho_pow_mul_sq_fourierCoeff}
\difficulty{4}
Let $n$ be a natural number, let $f:\{−1,1\}^n\to\mathbb R$, and let $0<\rho\le 1$.
Write $\widehat f(S)$ for the uniform average of $f(x)\prod_{i\in S}x_i$, and let
$W_1(f)=\sum_{|S|=1}\widehat f(S)^2$. If $\operatorname{Stab}_\rho(f)$ is the expected
product $f(X)f(Y)$ when $X$ is uniform on the Boolean cube and, independently in each
coordinate, $Y_i=X_i$ with probability $(1+\rho)/2$ and $Y_i=-X_i$ otherwise, then $\rho
W_1(f)\le \operatorname{Stab}_\rho(f)-(\mathbb E f)^2$.
\end{theorem}

\begin{theorem}[Level-one Fourier weight bound]
\lean{FABL.levelOneInequality_parameterized}
\label{FABL.levelOneInequality_parameterized}
\leanok
\uses{FABL.fourierWeightAtLevel, FABL.mean, FABL.rho_mul_fourierWeightAtLevel_one_le_stability_sub_mean_sq, FABL.smallSetExpansion}
\difficulty{3}
Let $f:\{-1,1\}^n\to\mathbb R$ take only the values $0$ and $1$, and let its mean be
$\alpha$. For every real number $\rho$ with $0<\rho\le 1$, the sum of the squared
Fourier coefficients of $f$ over all singleton coordinate sets is at most
$\rho^{-1}\left(\alpha^{2/(1+\rho)}-\alpha^2\right)$.
\end{theorem}

\begin{theorem}[Tail bound for low degree functions]
\lean{FABL.lowDegree_tail_bound_parameter}
\proofsource{Analysis-of-Boolean-Functions-by-Ryan-ODonnell}{
  pdf-ffbe9c7b3306-p263-b002,
  pdf-ffbe9c7b3306-p263-b003
}
\label{FABL.lowDegree_tail_bound_parameter}
\leanok
\uses{FABL.fourierDegree, FABL.lowDegree_uniformLpNorm_le, FABL.pmfEventProbability_congr, FABL.pmfEventProbability_le_expectation_div, FABL.pmfEventProbability_uniformPMF_eq_uniformProbability, FABL.pmfExpectation, FABL.pmfExpectation_uniformPMF_eq_expect, FABL.uniformLpNorm, FABL.uniformLpNorm_nonneg, FABL.uniformLpNorm_rpow, FABL.uniformPMF, FABL.uniformProbability}
\difficulty{6}
Let $f$ be a real-valued function on the $n$-dimensional Boolean cube with Fourier
degree at most $k$, where $k\in\mathbb N$. For real numbers $q,t$ with $2\le q$ and
$t>0$, if the uniform $L^2$ norm of $f$ is positive, then $\Pr[t\|f\|_2\le |f(x)|]\le
(q^{k/2}/t)^q$ for uniformly random $x$ in the cube.
\end{theorem}

\begin{theorem}[Low degree tail bound]
\lean{FABL.lowDegree_tail_bound}
\statementsource{Analysis-of-Boolean-Functions-by-Ryan-ODonnell}{
  pdf-ffbe9c7b3306-p263-b002,
  pdf-ffbe9c7b3306-p263-b003
}
\label{FABL.lowDegree_tail_bound}
\leanok
\uses{FABL.fourierDegree, FABL.lowDegree_tail_bound_parameter, FABL.uniformLpNorm, FABL.uniformProbability}
\difficulty{5}
% dataset-pipeline: stub -- prose pending build_informal_statements.py
Statement pending regeneration.
\end{theorem}
